\documentclass[11pt,a4paper]{article}

\usepackage{fontspec}
\usepackage{amsmath}
\usepackage{graphicx}
\usepackage{booktabs}
\usepackage{hyperref}

\setmainfont{Latin Modern Roman}

\title{Anatomy of a Synthetic Paper}
\author{Repository Fixture \and Second Author}
\date{2026-09-06}% ISO date uses hyphens, not dashes. chktex 8

\begin{document}
\maketitle

\begin{abstract}
  This fixture contains every structural element the semantic recovery
  engine must recognize in one paper: title block, abstract, numbered
  sections, a figure, a table, a display equation, a footnote, and a
  bibliography with in-text citations.
\end{abstract}

\section{Introduction}
Prior work on layout recovery~\cite{smith2020docs} and citation
resolution~\cite{doe2021refs} motivates this composite fixture. The
paragraph below is deliberately long so it wraps onto a second page and
exercises cross-page paragraph merging.
Synthetic filler sentences continue the same paragraph across the page
boundary, because the reading flow must emit continuation evidence for
semantic paragraph assembly rather than treating the fragments as
unrelated blocks. Additional filler lines follow so the break happens
inside this paragraph and not before the next section heading. This
paragraph keeps growing with more synthetic academic prose about
document compilers, evidence pipelines, and reading-order graphs.
Repetition is deliberate: it inflates vertical extent without
introducing new structural elements, isolating the cross-page
continuation case.

\subsection{Motivation}
A subsection heading must recover with level two, nested inside the
section tree. Inline math such as \(f(x) = x^{2}\) stays inside its
paragraph as a mark, not a node.

\section{Methods}
The pipeline recovers structure in stages. The placeholder drawn below
covers a region, and a synthetic footnote\footnote{The footnote body
  must be recovered and linked back to this marker.} adds reference
evidence. The linear identity anchors the numeric pipeline:
\begin{equation}
  y = 2x + 1
\end{equation}
\begin{figure}[htbp]
  \centering
  \rule{0.5\textwidth}{3cm}%
  \caption{Placeholder figure for composite recovery.}
\end{figure}

\section{Results}
The table reported below lists the synthetic metrics.

\begin{table}[htbp]
  \centering
  \begin{tabular}{lrr}
    \toprule
    Region type & Count & Confidence \\
    \midrule
    Paragraph   & 9     & 0.93       \\
    Heading     & 4     & 0.90       \\
    Figure      & 1     & 0.84       \\
    Table       & 1     & 0.80       \\
    \bottomrule
  \end{tabular}
  \caption{Composite recovery metrics.}
\end{table}

\section{Conclusion}
The conclusion closes the body before the bibliography. Recovery must
keep it a plain paragraph under the last numbered section.

\begin{thebibliography}{9}
  \bibitem{smith2020docs}
  A.~Smith.
  Synthetic references for composite regression tests.
  \emph{Fixture Journal}, 1(1):1--2, 2020.
  \bibitem{doe2021refs}
  J.~Doe.
  Citation marks in academic PDFs.
  \emph{Layout Letters}, 2(3):10--12, 2021.
\end{thebibliography}

\end{document}
